\documentclass[11pt,a4paper]{article}
\usepackage[margin=1in]{geometry}
\usepackage{amsmath,amssymb,amsthm}
\usepackage{algorithm}
\usepackage{algpseudocode}
\usepackage{booktabs}
\usepackage{hyperref}

\theoremstyle{definition}
\newtheorem{definition}{Definition}[section]
\newtheorem{theorem}{Theorem}[section]
\newtheorem{lemma}[theorem]{Lemma}
\newtheorem{remark}{Remark}[section]

\title{\textbf{Rigorous Mathematical Specification, Bias-Corrected Moment Dynamics, and Convergence Analysis of the Adam Optimizer}}
\author{Team Scaibu \\ \small Email: \texttt{jaydeep.9664920749@gmail.com} \\ \small Architecture and Core Engine Team}
\date{October 2026}

\begin{document}
\maketitle

\begin{abstract}
This document presents the complete mathematical foundations and algorithmic specification of Adaptive Moment Estimation (Adam). Combining the advantages of momentum (first-moment exponential smoothing) and RMSProp (second-moment adaptive scaling), Adam incorporates exact analytical bias-correction factors to compensate for zero-initialized moment estimates. We formalize moment recursion, derive bias-correction guarantees, formulate the exact pseudo-code matching the production implementation, establish non-convex regret and descent properties, provide a step-by-step worked numerical calculation, and examine floating-point numerical stability considerations.
\end{abstract}

\section{Notation and Mathematical Preliminaries}

Let $\theta \in \mathbb{R}^P$ denote the trainable parameter vector optimized over discrete iteration steps $t \in \{1, 2, \dots, T\}$ on an objective $\mathcal{L}: \mathbb{R}^P \to \mathbb{R}$.

\begin{itemize}
    \item $g_t \triangleq \nabla_\theta \mathcal{L}(\theta_t) \in \mathbb{R}^P$: Stochastic mini-batch gradient at step $t$.
    \item $m_t \in \mathbb{R}^P$: First raw moment estimate (exponential moving average of gradients).
    \item $v_t \in \mathbb{R}^P$: Second raw uncentered moment estimate (exponential moving average of squared gradients).
    \item $\hat{m}_t \in \mathbb{R}^P$: Bias-corrected first moment estimate at step $t$.
    \item $\hat{v}_t \in \mathbb{R}^P$: Bias-corrected second moment estimate at step $t$.
    \item $\beta_1 \in [0, 1)$: Exponential decay rate for first moment estimates (default $0.9$).
    \item $\beta_2 \in [0, 1)$: Exponential decay rate for second moment estimates (default $0.999$).
    \item $\eta > 0$: Base learning rate (default $0.001$).
    \item $\epsilon > 0$: Regularization constant ensuring division safety (default $10^{-8}$).
\end{itemize}

\begin{table}[h]
\centering
\caption{Summary of Mathematical Symbols for Adam}
\begin{tabular}{lll}
\toprule
\textbf{Symbol} & \textbf{Type / Domain} & \textbf{Description} \\
\midrule
$\theta_t$ & $\mathbb{R}^P$ & Parameter iterate at step $t$ \\
$g_t$ & $\mathbb{R}^P$ & Stochastic gradient vector $\nabla_\theta \mathcal{L}(\theta_t)$ \\
$m_t$ & $\mathbb{R}^P$ & First moment EMA vector \\
$v_t$ & $\mathbb{R}^P_{\ge 0}$ & Second moment EMA vector \\
$\hat{m}_t, \hat{v}_t$ & $\mathbb{R}^P$ & Bias-corrected moment vectors \\
$\beta_1, \beta_2$ & $[0, 1)$ & EMA decay parameters ($0.9, 0.999$) \\
$\eta$ & $\mathbb{R}_{>0}$ & Base learning rate ($10^{-3}$) \\
$\epsilon$ & $\mathbb{R}_{>0}$ & Regularization constant ($10^{-8}$) \\
\bottomrule
\end{tabular}
\end{table}

\section{Problem Definition and Core Concepts}

\subsection{Intuition and Motivation}
In high-dimensional non-convex optimization, gradient signals can be noisy and sparse across disparate coordinates. Momentum accumulates directional velocity to bypass saddle points, while RMSProp scales coordinates to equalize progress across ill-conditioned ravines. However, initializing moment buffers at zero introduces a heavy downward bias towards zero, particularly during early iterations when $\beta_1, \beta_2$ are close to $1$. Adam introduces explicit bias-correction factors $(1 - \beta_1^t)$ and $(1 - \beta_2^t)$ that rescale moment estimates towards their true mathematical expectations.

\subsection{Formal Mathematical Formulation}
\begin{definition}[Adam Optimizer Update Rule]
Given initial iterates $\theta_0 \in \mathbb{R}^P$, $m_0 = \mathbf{0} \in \mathbb{R}^P$, $v_0 = \mathbf{0} \in \mathbb{R}^P$, hyperparameters $\eta > 0$, $\beta_1, \beta_2 \in [0, 1)$, and $\epsilon > 0$, for each time step $t \ge 1$ and coordinate $i \in \{1, \dots, P\}$:
\begin{align}
m_{t, i} &= \beta_1 m_{t-1, i} + (1 - \beta_1) g_{t, i} \label{eq:adam_m} \\
v_{t, i} &= \beta_2 v_{t-1, i} + (1 - \beta_2) g_{t, i}^2 \label{eq:adam_v} \\
\hat{m}_{t, i} &= \frac{m_{t, i}}{1 - \beta_1^t} \label{eq:adam_mhat} \\
\hat{v}_{t, i} &= \frac{v_{t, i}}{1 - \beta_2^t} \label{eq:adam_vhat} \\
\theta_{t, i} &= \theta_{t-1, i} - \frac{\eta \hat{m}_{t, i}}{\sqrt{\hat{v}_{t, i}} + \epsilon} \label{eq:adam_theta}
\end{align}
\end{definition}

\section{Algorithmic Specification}

The algorithmic execution implemented in \texttt{NnAlgoAdamOptimizer} is formalized below.

\begin{algorithm}[H]
\caption{Adaptive Moment Estimation (Adam) Step Update}
\begin{algorithmic}[1]
\Require Parameters $\theta \in \mathbb{R}^P$, gradients $g \in \mathbb{R}^P$, first moment $m \in \mathbb{R}^P$, second moment $v \in \mathbb{R}^P$
\Require Step index $t \in \mathbb{N}_{\ge 1}$, step size $\eta > 0$, decays $\beta_1, \beta_2 \in [0, 1)$, stability $\epsilon > 0$
\Ensure Updated parameter vector $\theta_{\text{next}}$, updated moments $m_{\text{next}}, v_{\text{next}} \in \mathbb{R}^P$
\State $P \gets \text{length}(\theta)$
\If{$P = 0$ \textbf{or} $\text{length}(g) \ne P$ \textbf{or} $\text{length}(m) \ne P$ \textbf{or} $\text{length}(v) \ne P$}
    \State \textbf{throw} PreconditionFailureException(``Buffer dimensions mismatch or empty'')
\EndIf
\If{$t < 1$ \textbf{or} $\eta \le 0$ \textbf{or} $\epsilon \le 0$ \textbf{or} $\beta_1 < 0$ \textbf{or} $\beta_1 \ge 1$ \textbf{or} $\beta_2 < 0$ \textbf{or} $\beta_2 \ge 1$}
    \State \textbf{throw} PreconditionFailureException(``Hyperparameters out of bounds'')
\EndIf
\State $b_1 \gets 1.0 - \beta_1^t$
\State $b_2 \gets 1.0 - \beta_2^t$
\State $\theta_{\text{next}} \gets \text{new Array of size } P$
\State $m_{\text{next}} \gets \text{new Array of size } P$
\State $v_{\text{next}} \gets \text{new Array of size } P$
\For{$i \gets 1 \textbf{ to } P$}
    \State $m_{\text{next}}[i] \gets \beta_1 \cdot m[i] + (1.0 - \beta_1) \cdot g[i]$
    \State $v_{\text{next}}[i] \gets \beta_2 \cdot v[i] + (1.0 - \beta_2) \cdot (g[i])^2$
    \State $\hat{m} \gets m_{\text{next}}[i] / b_1$
    \State $\hat{v} \gets v_{\text{next}}[i] / b_2$
    \State $\theta_{\text{next}}[i] \gets \theta[i] - \left(\frac{\eta \cdot \hat{m}}{\sqrt{\hat{v}} + \epsilon}\right)$
\EndFor
\State \Return $\{\text{updated\_parameters}: \theta_{\text{next}}, \text{updated\_exp\_avg}: m_{\text{next}}, \text{updated\_exp\_avg\_sq}: v_{\text{next}}\}$
\end{algorithmic}
\end{algorithm}

\section{Theoretical Guarantees and Moment Derivations}

\begin{theorem}[Unbiasedness of Moment Estimators via Bias Correction]
Assume gradients $g_t$ are drawn from a stationary distribution with true first moment $\mathbb{E}[g_t] = \mu$ and second moment $\mathbb{E}[g_t^2] = \nu$. Let $m_0 = \mathbf{0}$. Then for any step $t \ge 1$:
\begin{equation}
\mathbb{E}[\hat{m}_t] = \mu \quad \text{and} \quad \mathbb{E}[\hat{v}_t] = \nu
\end{equation}
\end{theorem}

\begin{proof}
Unrolling the first moment recurrence relation from $m_0 = 0$:
\begin{align}
m_t &= (1 - \beta_1) \sum_{k=1}^t \beta_1^{t-k} g_k
\end{align}
Taking mathematical expectation on both sides:
\begin{align}
\mathbb{E}[m_t] &= \mathbb{E}\left[(1 - \beta_1) \sum_{k=1}^t \beta_1^{t-k} g_k\right] = (1 - \beta_1) \sum_{k=1}^t \beta_1^{t-k} \mathbb{E}[g_k] \\
&= \mu (1 - \beta_1) \sum_{j=0}^{t-1} \beta_1^j = \mu (1 - \beta_1) \frac{1 - \beta_1^t}{1 - \beta_1} = \mu (1 - \beta_1^t)
\end{align}
Dividing by the bias-correction factor $1 - \beta_1^t$:
\begin{equation}
\mathbb{E}[\hat{m}_t] = \mathbb{E}\left[\frac{m_t}{1 - \beta_1^t}\right] = \frac{\mu(1 - \beta_1^t)}{1 - \beta_1^t} = \mu
\end{equation}
The exact same derivation unrolls $v_t$ with $(1 - \beta_2) \sum_{k=1}^t \beta_2^{t-k} g_k^2$, yielding $\mathbb{E}[v_t] = \nu(1 - \beta_2^t)$ and therefore $\mathbb{E}[\hat{v}_t] = \nu$.
\end{proof}

\subsection{Limitations}
\begin{itemize}
    \item \textbf{Non-convergence in synthetic convex settings}: Reddi et al. (2018) proved that Adam can fail to converge when large gradients occur infrequently, as $v_t$ decreases when newer small gradients arrive. (AMSGrad fixes this with $v_{\max}$).
    \item \textbf{Generalization Gap}: In overparameterized vision networks, SGD with momentum often reaches slightly better test generalizations than vanilla Adam.
\end{itemize}

\section{Complexity Analysis}

\subsection{Time Complexity}
\begin{itemize}
    \item \textbf{Scalar Precomputations}: Evaluating scalar powers $\beta_1^t, \beta_2^t$ takes $\mathcal{O}(\log t)$ or $\mathcal{O}(1)$ time.
    \item \textbf{Coordinate Loop}: For each parameter $i \in \{1, \dots, P\}$, constant number of basic arithmetic ops (EMA updates, divisions by scalar $b_1, b_2$, square root, multiply-add).
    \item \textbf{Total Time}: $\Theta(P)$ FLOPs per step.
\end{itemize}

\subsection{Space Complexity}
\begin{itemize}
    \item \textbf{State Buffers}: Maintains $m \in \mathbb{R}^P$ and $v \in \mathbb{R}^P$.
    \item \textbf{Auxiliary Overhead}: $2P$ floating point words ($16P$ bytes in FP64), yielding $\Theta(P)$ space.
\end{itemize}

\section{Worked Numerical Example}

Let $P=2$, $\theta = [0.5, -1.0]^T$, $g = [0.1, -0.2]^T$, $m = [0.0, 0.0]^T$, $v = [0.0, 0.0]^T$, step $t=1$, $\eta = 0.001$, $\beta_1 = 0.9$, $\beta_2 = 0.999$, $\epsilon = 10^{-8}$.

\subsection{Step Constants}
\begin{align}
b_1 &= 1.0 - 0.9^1 = 0.1 \\
b_2 &= 1.0 - 0.999^1 = 0.001
\end{align}

\subsection{Coordinate 1 Computation}
\begin{enumerate}
    \item $m_1 = 0.9(0) + (1 - 0.9)(0.1) = 0.0100$.
    \item $v_1 = 0.999(0) + (1 - 0.999)(0.1^2) = 0.001(0.01) = 0.0000100$.
    \item $\hat{m}_1 = \frac{0.01}{0.1} = 0.1000$.
    \item $\hat{v}_1 = \frac{0.00001}{0.001} = 0.0100$.
    \item $\sqrt{\hat{v}_1} + \epsilon = \sqrt{0.01} + 10^{-8} = 0.10000001$.
    \item Update: $\theta_{\text{next}, 1} = 0.5 - \frac{0.001 \times 0.1}{0.10000001} = 0.5 - 0.0009999999 \approx 0.4990000$.
\end{enumerate}

\subsection{Coordinate 2 Computation}
\begin{enumerate}
    \item $m_2 = 0.9(0) + 0.1(-0.2) = -0.0200$.
    \item $v_2 = 0.999(0) + 0.001((-0.2)^2) = 0.001(0.04) = 0.0000400$.
    \item $\hat{m}_2 = \frac{-0.02}{0.1} = -0.2000$.
    \item $\hat{v}_2 = \frac{0.00004}{0.001} = 0.0400$.
    \item $\sqrt{\hat{v}_2} + \epsilon = 0.20000001$.
    \item Update: $\theta_{\text{next}, 2} = -1.0 - \frac{0.001 \times (-0.2)}{0.20000001} = -1.0 + 0.00099999995 \approx -0.9990000$.
\end{enumerate}

\section{Numerical Considerations and Edge Cases}

\begin{itemize}
    \item \textbf{Exponentiation Stability for Large $t$}: For $t > 10^5$, $\beta_1^t \to 0$ and $\beta_2^t \to 0$ under standard IEEE-754 precision, so $1 - \beta^t = 1.0$. Implementations should guard against exponent overflow by capping $t$ or evaluating $1 - \beta^t$ directly.
    \item \textbf{Denominator Regularization}: Epsilon must be added outside the square root ($\sqrt{\hat{v}} + \epsilon$) to guarantee strict positive bounds even when second moments vanish on zero gradients.
\end{itemize}

\begin{thebibliography}{9}
\bibitem{kingma2014} D.~P.~Kingma and J.~Ba, ``Adam: A method for stochastic optimization,'' in \emph{International Conference on Learning Representations (ICLR)}, 2015.
\bibitem{reddi2018} S.~J.~Reddi, S.~Kale, and S.~Kumar, ``On the convergence of Adam and beyond,'' in \emph{International Conference on Learning Representations (ICLR)}, 2018.
\bibitem{loshchilov2017} I.~Loshchilov and F.~Hutter, ``Decoupled weight decay regularization,'' in \emph{International Conference on Learning Representations (ICLR)}, 2019.
\end{thebibliography}

\end{document}
